\section{A quantitative bridge from model uncertainty to event records}
\label{sec:bridge}

This section constructs the missing marked-path comparison. It does not
deduce it from the older CDF comparison. The two comparisons have different
targets and remain distinct.

\subsection{Two-model input and the retained record}
Let $Q,\widehat Q$ be finite irreducible generators on the same encoded
state space, with invariant laws $\pi,\widehat\pi$ and the same nonempty
proper event set $B$. Use a common rational
$R\ge2\max\{\max_x[-Q(x,x)],\max_x[-\widehat Q(x,x)]\}$ and put
$P=I+Q/R$, $\widehat P=I+\widehat Q/R$.
Assume both normalized additive symmetrizations have gap at least $g>0$
and both invariant laws have mass floor $\zeta>0$.
Stationary reverses are denoted by stars. Common transition support is a
premise for the finite divergence envelopes below.

The event record has the same meaning as in the accepted first-event
construction: the quantized first time, terminal state and ordered changes
in the last $H$ normalized time units, truncated at the supplied initial
time. The comparison first proves a statement for continuous records;
applying any common finite timestamp quantization cannot increase TV.

Define finite bounds
\begin{align}
\chi_B&\ge D_{\rm KL}(\pi(\cdot\mid B)\Vert\widehat\pi(\cdot\mid B)),\\
\kappa_+&\ge\sup_xD_{\rm KL}(P(x,\cdot)\Vert\widehat P(x,\cdot)),\\
\kappa_-&\ge\sup_xD_{\rm KL}(P^*(x,\cdot)\Vert\widehat P^*(x,\cdot)),\\
d_p&\ge|\log(\pi(B)/\widehat\pi(B))|.
\end{align}
These are model-confidence bounds, not assumed observations of a rare
trajectory. A local experimental construction of them is given below.

\subsection{Square-root distance of marked proposals}
Write
$\mathsf h^2(\nu,\omega)=\int(\sqrt{d\nu}-\sqrt{d\omega})^2$
also for subprobability measures. For probabilities,
$\TV(\nu,\omega)\le\mathsf h(\nu,\omega)$ and
$\mathsf h^2(\nu,\omega)\le D_{\rm KL}(\nu\Vert\omega)$.

Use the accepted length-$L$ marked-Palm proposal: its initial-mark branch
has probability $1/(1+L)$ and its internal-mark branch has probability
$L/(1+L)$. The two models use the same branch, Poisson count, marked
position and uniform timestamps; only their conditional rare state and
forward/reverse updates differ.

\begin{lemma}[Proposal information bound]\label{lem:palmKL}
For the two marked proposals $\mathsf P_L,\widehat{\mathsf P}_L$,
\begin{equation}
D_{\rm KL}(\mathsf P_L\Vert\widehat{\mathsf P}_L)
\le D_L:=
\chi_B+\frac{L(L+2)}{2(L+1)}(\kappa_++\kappa_-).
\label{eq:DL}
\end{equation}
\end{lemma}
\begin{proof}
Retain the branch and marked position as auxiliary variables. The relative
entropy chain rule first pays $\chi_B$ at the marked state and then the
appropriate row divergences for all generated updates. The auxiliary
count/time laws are identical. In the internal-mark branch, conditional
on $n=M+1$ events, the marked position is uniform in $1,\ldots,n$.
Its expected backward count is $(n+1)/2$ and its expected forward count
is $(n-1)/2$. Since $\mathbb E M=L$, weighting the two branches gives
expected forward and backward counts each equal to
$L(L+2)/(2(L+1))$. Dropping auxiliary variables can only decrease KL.
\end{proof}

Let $G$ mean that the first hit within the block lies in $(H,L]$, and
$N_B$ count its potential-event states in $B$, including time zero and
null updates. Both proposals have the same measurable acceptance
function $w=\mathbf 1_G/N_B$. Let their acceptance means be
$\mu,\widehat\mu$.

\begin{lemma}[Accepted-submeasure comparison]\label{lem:acceptKL}
If $\mu,\widehat\mu\ge q_0>0$ and $z=\sqrt{D_L/q_0}$, then
\begin{equation}
\TV(\mathsf P_L(\cdot\mid\mathrm{accept}),
     \widehat{\mathsf P}_L(\cdot\mid\mathrm{accept}))\le z,
\qquad
|\log(\mu/\widehat\mu)|\le2z.
\label{eq:acceptKL}
\end{equation}
\end{lemma}
\begin{proof}
Multiplication by $0\le w\le1$ cannot increase square-root distance.
For the accepted submeasures $a=w\mathsf P_L$ and
$\widehat a=w\widehat{\mathsf P}_L$,
\begin{equation}
\mathsf h^2(a,\widehat a)=
(\sqrt\mu-\sqrt{\widehat\mu})^2+
\sqrt{\mu\widehat\mu}\,
\mathsf h^2(a/\mu,\widehat a/\widehat\mu).
\label{eq:submass}
\end{equation}
The left side is at most $D_L$. The second term gives the TV bound.
For the mass bound, use
$|\log\sqrt\mu-\log\sqrt{\widehat\mu}|
\le|\sqrt\mu-\sqrt{\widehat\mu}|/\sqrt{q_0}$.
\end{proof}

The acceptance denominator is $q_0$, not $\pi(B)$. Under the accepted
uniform block conditions it can be taken as $g/8$.
Because the stationary block success probability is
$q_G=p(1+L)\mu$, the two geometric parameters obey
\begin{equation}
|\log(q_G/\widehat q_G)|\le d_p+2z.
\end{equation}
The Bernoulli coupling bound therefore gives
\begin{equation}
\TV(\operatorname{Geom}_0(q_G),
    \operatorname{Geom}_0(\widehat q_G))
\le\min\{1,e^{d_p+2z}-1\}.
\label{eq:geom-model}
\end{equation}
The product of this waiting index and an independent accepted path
differs by at most the sum of \eqref{eq:acceptKL} and
\eqref{eq:geom-model}.

\subsection{Complete event-law comparison, including the initial state}
For an integer buffer $b$, let $\eta=e^{-gb}/\zeta<1$.
For each model separately let
\begin{equation}
S_Q=\frac{\eta}{1-\eta}
+\frac{1+\eta}{1-\eta}\frac{q_Q(H)+q_Q(b)}{q_{G,Q}}
+(1+\eta)q_Q(b),
\label{eq:SQ}
\end{equation}
and define $S_{\widehat Q}$ analogously. These are the full structural
bounds of the accepted block construction, including its joint
pre-block buffer charge and unnormalized initial survivor law.

\begin{theorem}[Joint first-event model stability]\label{thm:stability}
Suppose both models satisfy the block conditions and have acceptance
means at least $q_0$. For every common initial law $\nu$,
\begin{equation}
\TV(\Law_{\nu,Q}(\mathcal R_{H,s}),
    \Law_{\nu,\widehat Q}(\mathcal R_{H,s}))
\le
S_Q+S_{\widehat Q}
+\sqrt{b\kappa_+/2}
+z+\min\{1,e^{d_p+2z}-1\}.
\label{eq:master}
\end{equation}
An initial-law TV discrepancy adds its value. There is no inverse
stationary-rarity or mean-lifetime factor.
\end{theorem}
\begin{proof}
Compare each true event record with its own ideal block construction.
Those comparisons cost $S_Q$ and $S_{\widehat Q}$. Couple the observed
initial length-$b$ paths. Under Poisson uniformization their relative
entropy is at most $b\kappa_+$, since they have the same initial law;
Pinsker gives the displayed prefix cost. On a matching early hit both
algorithms return the same early record. On matching survival their late
outputs consist of the independent geometric index and successful path,
shifted by the same integer calendar offset. The preceding two lemmas
control that product law. All bounds are uniform in $\nu$. Coupling
initial states adds an initial-law error by contraction of a Markov
channel, and final quantization cannot increase TV.
\end{proof}

\subsection{Explicit envelopes acquired from local factors and rates}
Suppose $W,\widehat W$ are unnormalized Gibbs weights satisfying
\begin{equation}
|\log\widehat W(x)-\log W(x)|\le E
\quad\text{for every }x,
\qquad
|\log(\widehat Q(x,y)/Q(x,y))|\le a
\quad(x\ne y).
\label{eq:local-envelope}
\end{equation}
For the reversible true and learned local models in the accepted active
learner, both reverse kernels equal the corresponding forward kernel.

\begin{lemma}[Finite local-to-proposal bound]\label{lem:localKL}
Under these reversible premises and the common lazy uniformizer,
\begin{equation}
d_p=2E,\qquad
\chi_B=2E^2,\qquad
\kappa_+=\kappa_-=(e^a-1)^2
\label{eq:envelopes}
\end{equation}
are valid conservative bounds. In particular,
\begin{equation}
D_L\le 2E^2+(L+1)(e^a-1)^2.
\label{eq:reversibleDL}
\end{equation}
\end{lemma}
\begin{proof}
Both the full partition sum and its restriction to $B$ lie within
$e^{\pm E}$ of their learned counterparts, giving $d_p\le2E$.
The conditional density ratio on $B$ lies in $[e^{-2E},e^{2E}]$.
For any two probabilities with ratio in $[e^{-c},e^c]$,
convexity of minus the logarithm gives
$D_{\rm KL}\le c\tanh(c/2)\le c^2/2$; apply $c=2E$.

For a row, put $v=\sum_{y\ne x}\widehat P(x,y)\le1/2$.
The off-diagonal chi-square contribution relative to the learned row is
at most $(e^a-1)^2v$. The diagonal contribution is at most
$(e^a-1)^2v^2/(1-v)\le(e^a-1)^2v$.
Their sum is at most $(e^a-1)^2$. Relative entropy is bounded by
chi-square divergence. Reversibility supplies the same reverse bound,
and $L(L+2)/(L+1)<L+1$.
\end{proof}

For a nonreversible confidence family the same theorem applies with
separately certified reverse-row bounds. Stationary weight ratios and
incoming rates must construct them. The reversible active learner does
not silently acquire arbitrary driven kinetics.

